\documentclass[10pt,a4paper]{amsart}
\usepackage{amsmath,amssymb,amsthm,graphicx}
\usepackage{verbatim}
\newenvironment{reference}{\comment}{\endcomment}
%\newenvironment{reference}{}{}
\newenvironment{slogan}{\comment}{\endcomment}
\newenvironment{history}{\comment}{\endcomment}

% For commutative diagrams we use Xy-pic
\usepackage[all]{xy}

% We use 2cell for 2-commutative diagrams.
\xyoption{2cell}
% Theorem environments.
%
\theoremstyle{plain}
\newtheorem{theorem}[subsection]{Theorem}
\newtheorem{proposition}[subsection]{Proposition}
\newtheorem{lemma}[subsection]{Lemma}

\theoremstyle{definition}
\newtheorem{definition}[subsection]{Definition}
\newtheorem{example}[subsection]{Example}
\newtheorem{exercise}[subsection]{Exercise}
\newtheorem{situation}[subsection]{Situation}

\theoremstyle{remark}
\newtheorem{remark}[subsection]{Remark}
\newtheorem{remarks}[subsection]{Remarks}

\numberwithin{equation}{subsection}

% Macros
%
\def\lim{\mathop{\mathrm{lim}}\nolimits}
\def\colim{\mathop{\mathrm{colim}}\nolimits}
\def\Spec{\mathop{\mathrm{Spec}}}
\def\Hom{\mathop{\mathrm{Hom}}\nolimits}
\def\Ext{\mathop{\mathrm{Ext}}\nolimits}
\def\SheafHom{\mathop{\mathcal{H}\!\mathit{om}}\nolimits}
\def\SheafExt{\mathop{\mathcal{E}\!\mathit{xt}}\nolimits}
\def\Sch{\mathit{Sch}}
\def\Mor{\mathop{\mathrm{Mor}}\nolimits}
\def\Ob{\mathop{\mathrm{Ob}}\nolimits}
\def\Sh{\mathop{\mathit{Sh}}\nolimits}
\def\NL{\mathop{N\!L}\nolimits}
\def\CH{\mathop{\mathrm{CH}}\nolimits}
\def\proetale{{pro\text{-}\acute{e}tale}}
\def\etale{{\acute{e}tale}}
\def\QCoh{\mathit{QCoh}}
\def\Ker{\mathop{\mathrm{Ker}}}
\def\Im{\mathop{\mathrm{Im}}}
\def\Coker{\mathop{\mathrm{Coker}}}
\def\Coim{\mathop{\mathrm{Coim}}}

% Boxtimes
%
\DeclareMathSymbol{\boxtimes}{\mathbin}{AMSa}{"02}

%
% Macros for moduli stacks/spaces
%
\def\QCohstack{\mathcal{QC}\!\mathit{oh}}
\def\Cohstack{\mathcal{C}\!\mathit{oh}}
\def\Spacesstack{\mathcal{S}\!\mathit{paces}}
\def\Quotfunctor{\mathrm{Quot}}
\def\Hilbfunctor{\mathrm{Hilb}}
\def\Curvesstack{\mathcal{C}\!\mathit{urves}}
\def\Polarizedstack{\mathcal{P}\!\mathit{olarized}}
\def\Complexesstack{\mathcal{C}\!\mathit{omplexes}}
% \Pic is the operator that assigns to X its picard group, usage \Pic(X)
% \Picardstack_{X/B} denotes the Picard stack of X over B
% \Picardfunctor_{X/B} denotes the Picard functor of X over B
\def\Pic{\mathop{\mathrm{Pic}}\nolimits}
\def\Picardstack{\mathcal{P}\!\mathit{ic}}
\def\Picardfunctor{\mathrm{Pic}}
\def\Deformationcategory{\mathcal{D}\!\mathit{ef}}
\usepackage{cleveref}
\crefname{originallemma}{original lemma}{original lemmas}
\Crefname{originallemma}{Original lemma}{Original lemmas}
\makeatletter
\newlabel{lemma-integral-A}{{04DM}{source}{}{sourcepointer-1}{}}
\newlabel{lemma-integral-A@cref}{{[originallemma][0][]04DM}{[1][0][]source}{}{}{}}
\makeatother
\crefname{originallemma}{original lemma}{original lemmas}
\Crefname{originallemma}{Original lemma}{Original lemmas}
\makeatletter
\newlabel{lemma-integral-B}{{04DR}{source}{}{sourcepointer-2}{}}
\newlabel{lemma-integral-B@cref}{{[originallemma][0][]04DR}{[1][0][]source}{}{}{}}
\makeatother
\crefname{originallemma}{original lemma}{original lemmas}
\Crefname{originallemma}{Original lemma}{Original lemmas}
\makeatletter
\newlabel{lemma-integral-homeo-onto-image-C}{{04DW}{source}{}{sourcepointer-3}{}}
\newlabel{lemma-integral-homeo-onto-image-C@cref}{{[originallemma][0][]04DW}{[1][0][]source}{}{}{}}
\makeatother
\crefname{originallemma}{original lemma}{original lemmas}
\Crefname{originallemma}{Original lemma}{Original lemmas}
\makeatletter
\newlabel{lemma-property-C-closed-implies}{{04DV}{source}{}{sourcepointer-4}{}}
\newlabel{lemma-property-C-closed-implies@cref}{{[originallemma][0][]04DV}{[1][0][]source}{}{}{}}
\makeatother
\crefname{originallemma}{original lemma}{original lemmas}
\Crefname{originallemma}{Original lemma}{Original lemmas}
\makeatletter
\newlabel{lemma-property-C-implies}{{04DT}{source}{}{sourcepointer-5}{}}
\newlabel{lemma-property-C-implies@cref}{{[originallemma][0][]04DT}{[1][0][]source}{}{}{}}
\makeatother
\crefname{originallemma}{original lemma}{original lemmas}
\Crefname{originallemma}{Original lemma}{Original lemmas}
\makeatletter
\newlabel{sites-lemma-morphism-of-sites-almost-cocontinuous}{{04BA}{source}{}{sourcepointer-6}{}}
\newlabel{sites-lemma-morphism-of-sites-almost-cocontinuous@cref}{{[originallemma][0][]04BA}{[1][0][]source}{}{}{}}
\makeatother
\crefname{originallemma}{original lemma}{original lemmas}
\Crefname{originallemma}{Original lemma}{Original lemmas}
\makeatletter
\newlabel{lemma-property-B-implies}{{04DO}{source}{}{sourcepointer-7}{}}
\newlabel{lemma-property-B-implies@cref}{{[originallemma][0][]04DO}{[1][0][]source}{}{}{}}
\makeatother
\crefname{originallemma}{original lemma}{original lemmas}
\Crefname{originallemma}{Original lemma}{Original lemmas}
\makeatletter
\newlabel{sites-lemma-exactness-properties}{{04D6}{source}{}{sourcepointer-8}{}}
\newlabel{sites-lemma-exactness-properties@cref}{{[originallemma][0][]04D6}{[1][0][]source}{}{}{}}
\makeatother
\crefname{originallemma}{original lemma}{original lemmas}
\Crefname{originallemma}{Original lemma}{Original lemmas}
\makeatletter
\newlabel{sites-modules-lemma-exactness}{{04DB}{source}{}{sourcepointer-9}{}}
\newlabel{sites-modules-lemma-exactness@cref}{{[originallemma][0][]04DB}{[1][0][]source}{}{}{}}
\makeatother
\title{Consequences of integrality and universal injectivity for the small pushforward}
\author{The Stacks Project Authors}
\date{stacks-project:etale-cohomology.tex:local-snapshot-20261011; GNU FDL 1.2 or later}
\begin{document}
\maketitle
\noindent\textit{Editorial note.} All delivered mathematical text is the author's own, copied from the pinned Stacks Project chapter file etale-cohomology.tex; no statement or proof step was rewritten, reordered or removed.\par\medskip
\noindent\textit{Editorial note.} The proof is the author's complete argument for all seven clauses: it establishes the properties, derives clauses (1)-(3) from them by named implications and clauses (4)-(6) formally, and records that clause (7) is the almost cocontinuity already obtained.\par\medskip
\noindent\textit{Editorial note.} Every labelled reference occurring inside the counted spans is carried in assembly.external\_references as a non-load-bearing navigation pointer (same-chapter ones at their definition site, cross-chapter ones at the invocation inside the counted span), so no load-bearing mathematics is replaced by a stub. No bibliography citation occurs inside the counted spans.\par\medskip
\noindent\textit{Editorial note.} The checked chapter list for same-name collisions is the set of stacks chapters staged in this repository: chow, cohomology, divisors, etale-cohomology, flat, more-algebra, more-morphisms, morphisms, spaces-more-morphisms, varieties together with this card's chapter etale-cohomology.\par\medskip
\noindent\textit{Editorial note.} The preamble is the chapter's own native preamble excerpt already used for this source: the Xy-pic setup (native preamble lines 14-18) and the native theorem declarations with the mathematical macros (native preamble lines 154-217).\par\medskip
\noindent\textit{Allowed standard prerequisites (editorial).}\begin{itemize}
\item Properties (A), (B) and (C) of a morphism of schemes. Same chapter, the list of conditions and the results showing that an integral morphism and a universally injective integral morphism have the appropriate properties. Supplies the three properties established at the start of the proof.
\item Implications of the properties for the pushforward and for almost cocontinuity. Same chapter, the results deducing each clause of the statement from property (A), (B) or (C), together with the general results on exactness properties of morphisms of sites. Supplies the derivation of the seven clauses.
\end{itemize}\medskip
\ifdefined\hypertarget\hypertarget{sourcepointer-1}{}\fi
\noindent\textit{Source reference.} Original lemma 04DM: sources/stacks-etale-cohomology/pinned-original-source.tex, source lines 5852--5852. Navigation pointer invoked at proof line 6245; the counted span contains the pointer only and the cited result is not reproduced inside it, so the pointer is registered as a non-load-bearing navigation stub in the established form for Stacks material. No bibliography citation is involved.\par\medskip
\ifdefined\hypertarget\hypertarget{sourcepointer-2}{}\fi
\noindent\textit{Source reference.} Original lemma 04DR: sources/stacks-etale-cohomology/pinned-original-source.tex, source lines 6026--6026. Navigation pointer invoked at proof line 6188; the counted span contains the pointer only and the cited result is not reproduced inside it, so the pointer is registered as a non-load-bearing navigation stub in the established form for Stacks material. No bibliography citation is involved.\par\medskip
\ifdefined\hypertarget\hypertarget{sourcepointer-3}{}\fi
\noindent\textit{Source reference.} Original lemma 04DW: sources/stacks-etale-cohomology/pinned-original-source.tex, source lines 6162--6162. Navigation pointer invoked at proof line 6247; the counted span contains the pointer only and the cited result is not reproduced inside it, so the pointer is registered as a non-load-bearing navigation stub in the established form for Stacks material. No bibliography citation is involved.\par\medskip
\ifdefined\hypertarget\hypertarget{sourcepointer-4}{}\fi
\noindent\textit{Source reference.} Original lemma 04DV: sources/stacks-etale-cohomology/pinned-original-source.tex, source lines 6132--6132. Navigation pointer invoked at proof line 6250; the counted span contains the pointer only and the cited result is not reproduced inside it, so the pointer is registered as a non-load-bearing navigation stub in the established form for Stacks material. No bibliography citation is involved.\par\medskip
\ifdefined\hypertarget\hypertarget{sourcepointer-5}{}\fi
\noindent\textit{Source reference.} Original lemma 04DT: sources/stacks-etale-cohomology/pinned-original-source.tex, source lines 6104--6104. Navigation pointer invoked at proof line 6255; the counted span contains the pointer only and the cited result is not reproduced inside it, so the pointer is registered as a non-load-bearing navigation stub in the established form for Stacks material. No bibliography citation is involved.\par\medskip
\ifdefined\hypertarget\hypertarget{sourcepointer-6}{}\fi
\noindent\textit{Source reference.} Original lemma 04BA: sources/stacks-etale-cohomology/pinned-original-source.tex, source lines 6257--6257. Navigation pointer invoked at proof line 6257; the counted span contains the pointer only and the cited result is not reproduced inside it, so the pointer is registered as a non-load-bearing navigation stub in the established form for Stacks material. No bibliography citation is involved.\par\medskip
\ifdefined\hypertarget\hypertarget{sourcepointer-7}{}\fi
\noindent\textit{Source reference.} Original lemma 04DO: sources/stacks-etale-cohomology/pinned-original-source.tex, source lines 5920--5920. Navigation pointer invoked at proof line 6259; the counted span contains the pointer only and the cited result is not reproduced inside it, so the pointer is registered as a non-load-bearing navigation stub in the established form for Stacks material. No bibliography citation is involved.\par\medskip
\ifdefined\hypertarget\hypertarget{sourcepointer-8}{}\fi
\noindent\textit{Source reference.} Original lemma 04D6: sources/stacks-etale-cohomology/pinned-original-source.tex, source lines 6262--6262. Navigation pointer invoked at proof line 6262; the counted span contains the pointer only and the cited result is not reproduced inside it, so the pointer is registered as a non-load-bearing navigation stub in the established form for Stacks material. No bibliography citation is involved.\par\medskip
\ifdefined\hypertarget\hypertarget{sourcepointer-9}{}\fi
\noindent\textit{Source reference.} Original lemma 04DB: sources/stacks-etale-cohomology/pinned-original-source.tex, source lines 6264--6264. Navigation pointer invoked at proof line 6264; the counted span contains the pointer only and the cited result is not reproduced inside it, so the pointer is registered as a non-load-bearing navigation stub in the established form for Stacks material. No bibliography citation is involved.\par\medskip
\begin{lemma}
\label{lemma-integral-universally-injective}
Let $f : X \to Y$ be a morphism of schemes. Assume that $f$ is
universally injective and integral (for example a closed immersion).
Then
\begin{enumerate}
\item
$f_{small, *} :
\Sh(X_\etale)
\to
\Sh(Y_\etale)$
reflects injections and surjections,
\item
$f_{small, *} :
\Sh(X_\etale)
\to
\Sh(Y_\etale)$
commutes with pushouts and coequalizers (and more generally
finite connected colimits),
\item $f_{small, *}$ transforms surjections into surjections (on sheaves
of sets and on abelian sheaves),
\item the map
$f_{small}^{-1}f_{small, *}\mathcal{F} \to \mathcal{F}$
is surjective for any sheaf (of sets or of abelian groups)
$\mathcal{F}$ on $X_\etale$,
\item the functor $f_{small, *}$ is faithful (on sheaves of sets and
on abelian sheaves),
\item
$f_{small, *} :
\textit{Ab}(X_\etale)
\to
\textit{Ab}(Y_\etale)$
is exact, and
\item the functor
$Y_\etale \to X_\etale$, $V \mapsto X \times_Y V$ is
almost cocontinuous.
\end{enumerate}
\end{lemma}

\begin{proof}
By
Lemmas \ref{lemma-integral-A},
\ref{lemma-integral-B} and
\ref{lemma-integral-homeo-onto-image-C}
we know that the morphism $f$ has properties (A), (B), and (C).
Moreover, by
Lemma \ref{lemma-property-C-closed-implies}
we know that the functor $Y_\etale \to X_\etale$ is
almost cocontinuous. Now we have
\begin{enumerate}
\item property (C) implies (1) by
Lemma \ref{lemma-property-C-implies},
\item almost continuous implies (2) by
Sites, Lemma \ref{sites-lemma-morphism-of-sites-almost-cocontinuous},
\item property (B) implies (3) by
Lemma \ref{lemma-property-B-implies}.
\end{enumerate}
Properties (4), (5), and (6) follow formally from the first three, see
Sites, Lemma \ref{sites-lemma-exactness-properties}
and
Modules on Sites, Lemma \ref{sites-modules-lemma-exactness}.
Property (7) we saw above.
\end{proof}

\end{document}
